\documentclass[11pt,a4paper]{article}

\usepackage[T1]{fontenc}
\usepackage[utf8]{inputenc}
\usepackage[ngerman]{babel}
\usepackage[a4paper,top=2cm,bottom=2cm,left=2cm,right=2cm]{geometry}
\usepackage{helvet}
\renewcommand{\familydefault}{\sfdefault}
\usepackage{microtype}
\usepackage{enumitem}
\usepackage{titlesec}
\usepackage{graphicx}
\usepackage{xcolor}
\usepackage[hidelinks]{hyperref}

\definecolor{iugreen}{RGB}{0,169,160}

\setlength{\parindent}{0pt}
\setlength{\parskip}{3pt}
\titlespacing*{\section}{0pt}{7pt}{2pt}
\titlespacing*{\subsection}{0pt}{5pt}{1pt}
\titleformat{\section}{\normalfont\large\bfseries\color{iugreen}}{\thesection.}{0.5em}{}
\titleformat{\subsection}{\normalfont\normalsize\bfseries}{}{0em}{}
\setlist[itemize]{leftmargin=1.3em,topsep=1pt,itemsep=1pt,parsep=0pt,label=\textcolor{iugreen}{\textbullet}}
\setlist[enumerate]{leftmargin=1.5em,topsep=1pt,itemsep=1pt,parsep=0pt}

\linespread{1.04}

\newcommand{\lbl}[1]{\textbf{#1:}\;}

\begin{document}
\pagestyle{empty}

\begin{center}
{\Large\bfseries Cloud Programming (DLBSEPCP01\_D)  Aufgabe 1, Konzeptionsphase}\\[4pt]
{\small Florian Zimmermann \quad\textbar\quad Matrikelnummer 92012248}
\end{center}

\vspace{2pt}

\section{Szenario und Anforderungen}

Ein international tätiges Unternehmen benötigt eine öffentliche Website.
Besucher greifen weltweit und rund um die Uhr zu. Jeder Ausfall oder jede Verzögerung
bedeutet Reputationsschaden und Geschäftsverlust. Daraus folgt die Vorgabe einer weltweit
hochverfügbaren Bereitstellung.

\subsection{Funktionale Anforderungen}
\begin{itemize}
  \item Statisches Hosting einer HTML-Datei, öffentlich erreichbar
  \item Erreichbarkeit über HTTPS
  \item Globale Erreichbarkeit ohne spürbare Ladezeiten
\end{itemize}

\subsection{Nicht-funktionale Anforderungen}
\begin{itemize}
  \item \lbl{Hochverfügbarkeit} 24/7 Betrieb, Verfügbarkeit $\geq 99{,}9\,\%$
  \item \lbl{Automatische Skalierung} Skalierung bei Last, keine manuellen Eingriffe nötig 
  \item \lbl{Performance} minimale Ladezeiten
  \item \lbl{Sicherheit} Schutz vor DDoS- und Layer-7-Angriffen
  \item \lbl{Kosteneffizienz} serverlose / managed Dienste, Pay-per-Use
  \item \lbl{Reproduzierbarkeit} gesamte Infrastruktur als Infrastructure as Code (IaC)
\end{itemize}

\section{Architektur}

\subsection{Cloud-Anbieter: Microsoft Azure}
\begin{itemize}
  \item \lbl{AWS} sehr großes Angebot, aber Kreditkarte zur Registrierung nötig
    \item \lbl{Google Cloud} gutes Free Tier, aber weniger passendes Kursmaterial
   \item \lbl{Azure (gewählt)} Studentenguthaben (100 USD) ohne Kreditkarte, passender Service-Stack
\end{itemize}

\subsection{Verworfene Ansätze}
\begin{itemize}
  \item \lbl{Azure Storage + Front Door} für Student-Abos gesperrt
  \item \lbl{Static Web Apps allein} verfügbar und kostenlos, aber CDN ist eine Black Box
\end{itemize}

\subsection{Gewählte Lösung: Azure Static Web Apps + Cloudflare}
\begin{itemize}
  \item Azure Static Web Apps als serverlose Hosting-Plattform
  \item Cloudflare als vorgelagerte CDN- und Sicherheitsschicht:
        300+ PoPs, WAF, DDoS-Schutz, HTTPS, DNS der Custom Domain
\end{itemize}

\section{Komponenten}

\begin{itemize}
  \item \lbl{Azure Static Web Apps} hostet index.html, serverlos, hochverfügbar,
        öffentlicher HTTPS-Endpoint
  \item \lbl{Cloudflare} Edge davor: nächstgelegener Point of Presence (PoP), HTTPS-Terminierung, WAF/DDoS-Filter,
        Custom-Domain-Verwaltung
  \item \lbl{Terraform} verwaltet beide Provider (azurerm + cloudflare)
\end{itemize}
\vspace{2pt}
\enlargethispage{1.5cm}
\includegraphics[width=\linewidth,height=5cm,keepaspectratio]{c4.drawio.png}

\end{document}